\documentclass[10pt,a4paper]{amsart}
\usepackage[margin=19mm]{geometry}
\usepackage{amsmath,amssymb,mathtools}
\usepackage[scr=rsfs,cal=euler]{mathalfa}
\usepackage{xfrac}
\usepackage[unicode,hidelinks,bookmarks=false]{hyperref}
\newcommand{\ie}{i.e.\ }
\newcommand{\cf}{cf.\ }
\newcommand{\resp}{respectively\ }
\newcommand{\Subsection}{\subsection}
\providecommand{\nobreakdash}{\nobreak\hbox{-}}
\setlength{\emergencystretch}{2em}
\multlinegap0pt
\usepackage{bm}
\usepackage{bbm}
\usepackage{dsfont}
\usepackage{xspace}
\usepackage{cancel}
\usepackage{mathtools}
\usepackage{amsthm}
\makeatletter
\@ifundefined{lemma}{\newtheorem{lemma}{Lemma}}{}
\@ifundefined{proposition}{\newtheorem{proposition}[lemma]{Proposition}}{}
\makeatother
\makeatletter
\newcommand{\Omegabar}{\overline{\Omega}}
\newcommand{\R}{\mathbb{R}}
\def\norm#1{\|#1\|}
\expandafter\ifx\csname proof\endcsname\relax
\renewenvironment{proof}[1][\proofname]{\medskip \noindent {\bfseries #1. }}{\qed \bigskip}
\fi
\makeatother
\title[Quantitative and exact concavity principles for parabolic an]{Quantitative and exact concavity principles for parabolic and elliptic equations}
\author{Marco Gallo, Riccardo Moraschi, Marco Squassina}
\date{Source: arXiv:2504.09494v2 (CC BY 4.0)}
\begin{document}
\maketitle

\noindent\emph{Source and licence.} Quantitative and exact concavity principles for parabolic and elliptic equations. Version arXiv:2504.09494v2 (CC BY 4.0), distributed under the Creative Commons Attribution 4.0 International licence, as recorded in the arXiv licence metadata for this exact version and shown on the abstract page https://arxiv.org/abs/2504.09494v2. Full rights notice: licenses/arxiv-2504.09494-CC-BY-4.0.txt.\par\smallskip

\noindent\textit{Editorial scope.} Counted unit: the Lemma whose source label is \texttt{lem\_behav\_boundar\_omega} in the article (source0.tex lines 1457--1477, source label \texttt{lem\_behav\_boundar\_omega}), delivered together with the article's own complete original proof of it (source0.tex lines 1479--1527, one \texttt{proof} environment of 49 lines). No mathematics is written, repaired or re-derived: statement and proof are the article's own text line for line. Required context retained uncounted: eq\_ipost\_domain (lines 1429--1431); regularity conditions (lines 1433--1435); eq\_general\_bxu (lines 1437--1439); boundary and initial conditions (lines 1441--1447); eq\_sub\_super\_sol (lines 2062--2068); comparison principle dickstein (lines 2071--2080); Hopf lemma in parabolic case (lines 2176--2187). Every definition, display and label of those lines is retained unchanged. Omitted from this package and not redistributed: 1--1428, 1432--1432, 1436--1436, 1440--1440, 1448--1456, 1478--1478, 1528--2061, 2069--2070, 2081--2175, 2188--2314. Nothing inside a retained excerpt is omitted or altered: every retained body is delivered exactly as the article's lines are written, so no omission receipt of any kind accompanies this package. Citations: the retained excerpts carry no citation command at all, so no bibliography entry of the article is transcribed and no reference is redistributed. Reference adaptation: none was needed and none was made; every reference inside the delivered unit resolves inside it, so no unresolved reference or citation is produced. Printed slips kept literal: no printed word, symbol, hypothesis or number of the counted statement or of its proof was altered. Layout and class: the article's own class is \texttt{amsart} with packages including amsmath, amssymb, latexsym, esint, mathrsfs, mathtools, amsthm, verbatim, wasysym, geometry, babel, fontenc, setspace, appendix; the wrapper is the lane's pinned \texttt{amsart} editorial wrapper at 10pt on A4, which restates the theorem-like environments the retained text needs and adds the article's own macro definitions that the retained text needs (\texttt{\textbackslash Omegabar}, \texttt{\textbackslash R}, \texttt{\textbackslash norm}), with extra wrapper packages (bm, bbm, dsfont, xspace, cancel, mathtools). The delivered numbering is the unit's own. Assets: no retained span contains a figure environment, a graphics-inclusion command, a PDF-page inclusion, a table or a TikZ picture; no image file is copied into this package, the package declares no asset dependency and no image file is redistributed with it. Licence: version arXiv:2504.09494v2 of this article is distributed under the Creative Commons Attribution 4.0 International licence, as recorded in the arXiv licence metadata for this exact version and shown on the abstract page https://arxiv.org/abs/2504.09494v2. A full rights notice is delivered at licenses/arxiv-2504.09494-CC-BY-4.0.txt. Characters: every retained excerpt is pure ASCII, so the delivered engine, XeLaTeX with its default Latin Modern text font, reports no missing character.
	\begin{equation}\label{eq_ipost_domain}
	\hbox{$\Omega \subset \R^n$, $n\geq 2$, bounded convex domain which satisfies the interior sphere property.}
	\end{equation}

	\begin{align}\label{regularity conditions}
	u\in C^2_x(\Omega)\cap C^1_x(\Omegabar)\cap C^1_t((0,+\infty))\cap C(\Omegabar\times[0,+\infty)).
	\end{align} 

	\begin{equation}\label{eq_general_bxu}
	u_t-\Delta u=b(x,u, t) \quad\hbox{in }\Omega\times (0,+\infty).
	\end{equation}

	\begin{equation}\label{boundary and initial conditions}
		\begin{cases}
				u>0 & \text{in } \Omega\times (0,+\infty),\\
					u=0 & \text{on }\partial \Omega\times [0,+\infty),	 \\
					u =0 & \text{on } \Omegabar \times \{0\}. 
		\end{cases}
	\end{equation}

\begin{equation} \label{eq_sub_super_sol}
		\begin{cases}
			w_t-\Delta w = b(x,w,t)&\hbox{in }\Omega\times(0,T),\\
			w=0 &\hbox{on } \partial \Omega\times(0,T),\\
			w=u_0&\hbox{on } \Omega \times \{0\},
		\end{cases}
	\end{equation}

\begin{proposition}[Comparison and uniqueness, I]
\label{comparison principle dickstein} 
Let $\Omega \subset \R^n$ be a bounded open domain with the interior sphere property. 
Let $T>0$ and $b\colon\Omega\times(0,+\infty) \times (0,T)\to [0,+\infty)$ satisfying \hyperlink{H2}{\( (H2) \)} and assume $b(x,\cdot,t)$ continuous for all $(x,t)\in \Omega\times(0,T)$.
	Let $u_0\in C(\Omegabar)$ with $u_0\geq \delta d_\Omega$ for some $\delta>0$ or $u_0\equiv 0$. 
	Let $u,v\in C^2_x(\Omega)\cap C^1_t((0,T])\cap C(\Omegabar\times[0,T])$ and $(u-v)^+\in H^1(\Omega)$.
	Assume $u$ is a positive supersolution and $v$ is a subsolution of \eqref{eq_sub_super_sol}.
Then $u(\cdot,t)\geq v(\cdot,t)$ for all $0\leq t\leq T$. 
In particular, there exists at most one positive solution of \eqref{eq_sub_super_sol}.
\end{proposition}

\begin{lemma}[Parabolic Hopf lemma]
\label{Hopf lemma in parabolic case} 	
Let $\Omega$ be an open bounded set in $\R^n$ satisfying the interior sphere condition, $T>0$ and $u\in C^2_x(\Omega)\cap C^1_t((0,T])\cap C(\Omegabar\times[0,T])$.
	Suppose further
	\[ u_t+Lu\geq 0 \quad \text{in }\Omega\times(0,T]\]
	where \[Lu:=-\sum\limits_{i,j=1}^{n}a_{ij} D^2_{ij}u+\sum\limits_{i=1}^{n}b^iD_{i}u\] 
	with $a_{ij}, b_i \in L^{\infty}(\Omega\times(0,T])\cap C(\Omega\times(0,T])$. 
	Moreover, suppose that there exists a point $(\bar{x},\bar{t})\in \partial_p\Omega_T$ such that
	\[ u(\bar{x},\bar{t})<u(x,t)\quad \forall (x,t)\in C_r(\bar{x},\bar{t}). \]
	Let $\mathbf{e}\in \R^{n+1}$ be a direction such that $(\bar{x},\bar{t})+s\mathbf{e}\in \overline{C_r(\overline{x},\overline{t})}$ for some $s>0$, and assume that $\partial_{\mathbf{e}} u (\bar{x},\bar{t})$ exists. 
	Then $\partial_{\mathbf{e}} u (\bar{x},\bar{t})>0$.
\end{lemma}

\begin{lemma}[Boundary estimates]
\label{lem_behav_boundar_omega}
\label{lem_behav_t=0}
Let $\Omega$ satisfy \eqref{eq_ipost_domain}, and $u$ be a function such that \eqref{regularity conditions}, \eqref{eq_general_bxu} and \eqref{boundary and initial conditions} hold.
Suppose 	
$b\colon\Omega\times(0,+\infty) \times (0,+\infty)\to \R$ satisfies \hyperlink{H1}{\( (H1) \)}, and consider $M, T, q,\gamma$ therein; we can assume $T$ small enough in such a way $\norm{u}_{L^{\infty}(\Omega\times[0,T])}<M$.
 The following facts hold.

\begin{itemize}

\item[(i)] $\,$ \emph{[}Dealing with $ \Omega \times \{0\}$\emph{]}:
let $\varphi_1$ be the first Dirichlet eigenfunction on $\Omega$ with eigenvalue $\lambda_1$, normalized by $\norm{\varphi_1}_{\infty}=1$. 
Then, $$u(x,t) \geq C e^{-\lambda_1 t} t^{\frac{1+\gamma}{1-q}}\varphi_1(x) \quad \hbox{in $\Omega \times (0,T)$}$$
for some explicit $C=C(\norm{u}_{L^{\infty}(\Omega \times [0,T])}, q, \gamma, T)>0$. 

\item[(ii)] $\,$\emph{[}Dealing with $\partial \Omega \times \{0\}$\emph{]}:
for any $x_0 \in \partial \Omega$ and $\beta \in (0,2]$, there exists $\delta=\delta(x_0,T, \beta)>0$ such that
$$u(x_0+t\nu,t^{\beta})\geq \delta t^{\frac{2\beta(1+\gamma)+(2-\beta)(1-q) }{2(1-q)}} \quad \hbox{for $t\in (0,T)$ small},$$
where $\nu$ is the interior normal to $\partial \Omega$ in $x_0$. 
\end{itemize}
\end{lemma}

\begin{proof}
From \hyperlink{H1}{\( (H1) \)} with $M=\norm{u}_{L^{\infty}(\Omega \times [0,T])}$, we know that there exists $k=k(M, T)>0$ such that $u$ satisfies \eqref{boundary and initial conditions} and 
$$u_t-\Delta u=b(x,u,t)\geq k t^{\gamma} u^q \quad \text{in }\Omega\times (0, T).$$
Set 
$$w(x,t):= C e^{-\lambda_1 t} t^{\frac{1+\gamma}{1-q}} \varphi_1(x)$$
and $C := (\frac{1-q}{1+\gamma} k)^{\frac{1}{1-q}} 
$. 
A straightforward computation shows that $w$ satisfies \eqref{boundary and initial conditions} 
together with
$$w_t-\Delta w = C \frac{1+\gamma}{1-q} e^{-\lambda_1 t} t^{\frac{\gamma+q}{1-q}} \varphi_1(x) \quad \text{in }\Omega\times (0,+\infty)$$
and hence, since $q <1$ and by definition of $C$, 
$$w_t-\Delta w = 
\frac{1}{k} C^{1-q} \frac{1+\gamma}{1-q} e^{-(1-q) \lambda_1 t} \varphi_1^{1-q}(x) (k t^{\gamma} w^q) \leq k t^{\gamma} w^q \quad \text{in }\Omega\times (0,+\infty).$$
Then by Proposition \ref{comparison principle dickstein}, $u\geq w$ in $\Omega\times (0,T)$, which is the first claim.

\medskip

To show the second claim, assume first that $x_0 = 0 \in \partial \Omega$.
Let $\eta \in (0,\min\{1,T^{\beta}\})$ and define $w_\eta$ and $\Omega_\eta$ by 
	\[ w_\eta(x,t) :=\eta^{\frac{1 +\gamma}{1-q}}w(\eta^{-1/2}x,\eta^{-1}t), \quad \Omega_\eta:= \sqrt{\eta} \Omega \]
	which satisfy		
	\begin{equation*}
		\begin{cases}
			\partial_t w_\eta-\Delta w_\eta \leq 
			k t^{\gamma}(w_\eta)^q &\text{in }\Omega_\eta\times (0,+\infty),\\
			w_\eta>0 & \text{in } \Omega_\eta\times (0,+\infty), \\
			w_\eta=0 & \text{on }\partial \Omega_\eta\times (0,+\infty),\\ 
			w_\eta=0& \text{on } \Omega_\eta \times \{0\}.	
			\end{cases}
	\end{equation*}
Noticed that $\Omega_{\eta} \subset \Omega$, by Proposition \ref{comparison principle dickstein} we obtain $u\geq w_\eta$ in $\Omegabar_\eta\times (0,T)$. 
Now let $t=\eta^{\frac{1}{\beta}} \in (0, T)$ and, observed $0 \in \partial \Omega_{\eta}$, we choose $x=t\nu \in \Omega_{\eta}$, so that
\begin{equation}\label{eq_dim_wtnutbeta}
 u(t\nu,t^{\beta})\geq t^{\frac{\beta (1+\gamma)
	}{1-q}}w(t^{\frac{2-\beta}{2}}\nu,1). 
	\end{equation}
	If $\beta=2$ we obtain $u(t\nu,t^{2})\geq t^{\frac{2 (1+\gamma)
	}{1-q}}w(\nu,1)$, and the claim follows choosing $\delta :=w(\nu,1)$. If $\beta \in (0,2)$, 
by Lagrange theorem there exists $\xi_{t}\in (0,t^{\frac{2-\beta}{2}}\nu)$ such that (recall $w(0,1)=0$)
	$$w(t^{\frac{2-\beta}{2}}\nu,1) = t^{\frac{2-\beta}{2}}\nu\cdot D w(\xi_t,1).
	$$
	Being $\delta:= \frac{1}{2}\nu\cdot D w(0,1)>0$ thanks to  Lemma \ref{Hopf lemma in parabolic case}, for $t$ small enough we have $\nu\cdot D w(\xi_t,1)>\delta$ and hence we obtain
	\begin{equation}\label{eq_dim_estim_w}
	 w(t^{\frac{2-\beta}{2}}\nu,1) \geq \delta t^{\frac{2-\beta}{2}}\quad \hbox{for $t$ small}.
	 \end{equation}
	Thus, by \eqref{eq_dim_wtnutbeta} and \eqref{eq_dim_estim_w} we obtain $u(t\nu,t^{\beta})\geq \delta t^{\frac{2\beta(1+\gamma)+(2-\beta)(1-q) }{2(1-q)}}$.

The case of general $x_0$ can be obtained straightforwardly, noticed the equation related to $u \mapsto k t^{\gamma} u^q$ is $x$-translation invariant.
\end{proof}

\end{document}
